\documentclass[10pt]{article}

\usepackage[
    letterpaper,
    top=0.55in,
    bottom=0.60in,
    left=0.72in,
    right=0.72in
]{geometry}

\usepackage{amsmath,amssymb,amsfonts}
\usepackage{enumitem}
\usepackage{fancyhdr}
\usepackage{times}
\usepackage{mathtools}

% ------------------------------------------------------------
% Header
% ------------------------------------------------------------

\pagestyle{fancy}
\fancyhf{}
\fancyhead[L]{\small\itshape CS229 Problem Set \#1}
\fancyhead[R]{\small\thepage}
\renewcommand{\headrulewidth}{0pt}
\setlength{\headheight}{14pt}

% ------------------------------------------------------------
% Compact spacing
% ------------------------------------------------------------

\setlength{\parindent}{0pt}
\setlength{\parskip}{0pt}

\setlist[enumerate,1]{
    leftmargin=1.55em,
    itemsep=0.55em,
    topsep=0.35em,
    parsep=0pt,
    partopsep=0pt
}

\setlist[enumerate,2]{
    leftmargin=2.0em,
    itemsep=0.45em,
    topsep=0.30em,
    parsep=0pt,
    partopsep=0pt
}

\setlist[enumerate,3]{
    leftmargin=2.2em,
    itemsep=0.35em,
    topsep=0.25em,
    parsep=0pt,
    partopsep=0pt
}

\abovedisplayskip=5pt
\belowdisplayskip=5pt
\abovedisplayshortskip=4pt
\belowdisplayshortskip=4pt
\setlength{\jot}{2pt}

\newcommand{\ind}{\mathbf{1}}

\begin{document}

% ============================================================
% Title
% ============================================================

{\Large CS 229, Fall 2018}\\[-1pt]
{\LARGE Problem Set \#1: Supervised Learning}

\vspace{0.6em}
\hrule
\vspace{0.8em}

\begin{enumerate}

% ============================================================
% Problem 1
% ============================================================

\item \textbf{[40 points] Linear Classifiers (logistic regression and GDA)}

In this problem, we cover two probabilistic linear classifiers we have
covered in class so far. First, a discriminative linear classifier:
logistic regression. Second, a generative linear classifier: Gaussian
discriminant analysis (GDA). Both algorithms find a linear decision
boundary that separates the data into two classes, but make different
assumptions. Our goal in this problem is to get a deeper understanding
of the similarities and differences (and strengths and weaknesses) of
these two algorithms.

For this problem, we will consider two datasets, provided in the
following files:
\begin{enumerate}[label=\roman*.]
    \item \texttt{data/ds1\_\{train,valid\}.csv}
    \item \texttt{data/ds2\_\{train,valid\}.csv}
\end{enumerate}

Each file contains $m$ examples, one example
$(x^{(i)},y^{(i)})$ per row. In particular, the $i$th row contains
columns $x_0^{(i)}\in\mathbb{R}$,
$x_1^{(i)}\in\mathbb{R}$, and $y^{(i)}\in\{0,1\}$.
In the subproblems that follow, we will investigate using logistic
regression and Gaussian discriminant analysis (GDA) to perform binary
classification on these two datasets.

\begin{enumerate}[label=(\alph*)]

% ------------------------------------------------------------
% 1(a)
% ------------------------------------------------------------

\item \textbf{[10 points]}
In lecture we saw the average empirical loss for logistic regression:
\[
J(\theta)
=
-\frac{1}{m}
\sum_{i=1}^{m}
\left[
y^{(i)}
\log\left(h_\theta(x^{(i)})\right)
+
(1-y^{(i)})
\log\left(1-h_\theta(x^{(i)})\right)
\right],
\]
where
\[
y^{(i)}\in\{0,1\},
\qquad
h_\theta(x)=g(\theta^Tx),
\qquad
g(z)=\frac{1}{1+e^{-z}}.
\]

Find the Hessian $H$ of this function, and show that for any vector
$z$, it holds true that
\[
z^THz\geq 0.
\]

\textit{Hint:}
You may want to start by showing that
\[
\sum_i\sum_j z_i x_i x_j z_j
=
(x^Tz)^2
\geq 0.
\]
Recall also that
\[
g'(z)=g(z)(1-g(z)).
\]

\textit{Remark:}
This is one of the standard ways of showing that the matrix $H$ is
positive semi-definite, written $H\succeq0$.
This implies that $J$ is convex, and has no local minima other than
the global one. If you have some other way of showing $H\succeq0$,
you are also welcome to use your method instead of the one above.

% Solution to 1(a) will go here.
\textbf{Solution.}

Let
\[
h_i
:=
h_\theta(x^{(i)})
=
g(\theta^T x^{(i)}).
\]
For the $i$th training example, define
\[
\ell_i(\theta)
=
-
\left[
y^{(i)}\log h_i
+
(1-y^{(i)})\log(1-h_i)
\right].
\]
Then
\[
J(\theta)
=
\frac{1}{m}
\sum_{i=1}^m
\ell_i(\theta).
\]

Since
\[
g'(t)=g(t)(1-g(t)),
\]
we have
\[
\frac{\partial h_i}{\partial\theta_j}
=
h_i(1-h_i)x_j^{(i)}.
\]

Therefore,
\[
\begin{aligned}
\frac{\partial \ell_i}{\partial\theta_j}
&=
-\frac{y^{(i)}}{h_i}
\frac{\partial h_i}{\partial\theta_j}
+
\frac{1-y^{(i)}}{1-h_i}
\frac{\partial h_i}{\partial\theta_j}
\\
&=
-y^{(i)}(1-h_i)x_j^{(i)}
+
(1-y^{(i)})h_i x_j^{(i)}
\\
&=
(h_i-y^{(i)})x_j^{(i)}.
\end{aligned}
\]

Hence,
\[
\frac{\partial J}{\partial\theta_j}
=
\frac1m
\sum_{i=1}^m
(h_i-y^{(i)})x_j^{(i)},
\]
and thus
\[
\nabla J(\theta)
=
\frac1m
\sum_{i=1}^m
(h_i-y^{(i)})x^{(i)}.
\]

Now consider the $(j,k)$ entry of the Hessian:
\[
\begin{aligned}
H_{jk}
&=
\frac{\partial^2J}
{\partial\theta_j\partial\theta_k}
\\
&=
\frac1m
\sum_{i=1}^m
\frac{\partial h_i}{\partial\theta_k}
x_j^{(i)}
\\
&=
\frac1m
\sum_{i=1}^m
h_i(1-h_i)
x_j^{(i)}x_k^{(i)}.
\end{aligned}
\]

Therefore,
\[
\boxed{
H
=
\frac1m
\sum_{i=1}^m
h_i(1-h_i)
x^{(i)}(x^{(i)})^T
}.
\]

To show that $H$ is positive semidefinite, let
$z\in\mathbb{R}^n$ be arbitrary. Then
\[
\begin{aligned}
z^THz
&=
\frac1m
\sum_{i=1}^m
h_i(1-h_i)
z^Tx^{(i)}(x^{(i)})^Tz
\\
&=
\frac1m
\sum_{i=1}^m
h_i(1-h_i)
\left((x^{(i)})^Tz\right)^2.
\end{aligned}
\]

Since
\[
0<h_i<1,
\]
we have
\[
h_i(1-h_i)\ge0,
\]
and clearly
\[
\left((x^{(i)})^Tz\right)^2\ge0.
\]
Thus every term in the summation is nonnegative, so
\[
\boxed{
z^THz\ge0
}
\]
for every $z\in\mathbb{R}^n$.

Hence,
\[
\boxed{
H\succeq0
}.
\]

% ------------------------------------------------------------
% 1(b)
% ------------------------------------------------------------

\item \textbf{[5 points] Coding problem.}

Follow the instructions in \texttt{src/p01b\_logreg.py} to train a
logistic regression classifier using Newton's Method.

Starting with
\[
\theta=\vec{0},
\]
run Newton's Method until the updates to $\theta$ are small.
Specifically, train until the first iteration $k$ such that
\[
\|\theta_k-\theta_{k-1}\|_1<\epsilon,
\qquad
\epsilon=1\times10^{-5}.
\]

Make sure to write your model's predictions to the file specified in
the code.

% ------------------------------------------------------------
% 1(c)
% ------------------------------------------------------------

\item \textbf{[5 points]}

Recall that in GDA we model the joint distribution of $(x,y)$ by the
following equations:
\[
p(y)
=
\begin{cases}
\phi, & y=1,\\
1-\phi, & y=0,
\end{cases}
\]
\[
p(x\mid y=0)
=
\frac{1}
{(2\pi)^{n/2}|\Sigma|^{1/2}}
\exp\left(
-\frac12(x-\mu_0)^T\Sigma^{-1}(x-\mu_0)
\right),
\]
and
\[
p(x\mid y=1)
=
\frac{1}
{(2\pi)^{n/2}|\Sigma|^{1/2}}
\exp\left(
-\frac12(x-\mu_1)^T\Sigma^{-1}(x-\mu_1)
\right),
\]
where $\phi,\mu_0,\mu_1,$ and $\Sigma$ are the parameters of our model.

Suppose we have already fit $\phi,\mu_0,\mu_1,$ and $\Sigma$, and now
want to predict $y$ given a new point $x$.

To show that GDA results in a classifier that has a linear decision
boundary, show the posterior distribution can be written as
\[
p(y=1\mid x;\phi,\mu_0,\mu_1,\Sigma)
=
\frac{1}
{1+\exp\left(-(\theta^Tx+\theta_0)\right)},
\]
where
\[
\theta\in\mathbb{R}^n,
\qquad
\theta_0\in\mathbb{R}
\]
are appropriate functions of $\phi,\Sigma,\mu_0,$ and $\mu_1$.

% Solution to 1(c) will go here.

% ------------------------------------------------------------
% 1(d)
% ------------------------------------------------------------

\item \textbf{[7 points]}

For this part of the problem only, you may assume $n$ (the dimension
of $x$) is $1$, so that
\[
\Sigma=[\sigma^2]
\]
is just a real number, and likewise the determinant of $\Sigma$ is
given by
\[
|\Sigma|=\sigma^2.
\]

Given the dataset, we claim that the maximum likelihood estimates of
the parameters are given by
\[
\phi
=
\frac{1}{m}
\sum_{i=1}^{m}
\ind\{y^{(i)}=1\},
\]
\[
\mu_0
=
\frac{
\sum_{i=1}^{m}
\ind\{y^{(i)}=0\}x^{(i)}
}{
\sum_{i=1}^{m}
\ind\{y^{(i)}=0\}
},
\]
\[
\mu_1
=
\frac{
\sum_{i=1}^{m}
\ind\{y^{(i)}=1\}x^{(i)}
}{
\sum_{i=1}^{m}
\ind\{y^{(i)}=1\}
},
\]
and
\[
\Sigma
=
\frac{1}{m}
\sum_{i=1}^{m}
\left(
x^{(i)}-\mu_{y^{(i)}}
\right)
\left(
x^{(i)}-\mu_{y^{(i)}}
\right)^T.
\]

The log-likelihood of the data is
\[
\ell(\phi,\mu_0,\mu_1,\Sigma)
=
\log
\prod_{i=1}^{m}
p(x^{(i)},y^{(i)};\phi,\mu_0,\mu_1,\Sigma),
\]
which can be written as
\[
\ell(\phi,\mu_0,\mu_1,\Sigma)
=
\log
\prod_{i=1}^{m}
p(x^{(i)}\mid y^{(i)};\mu_0,\mu_1,\Sigma)
p(y^{(i)};\phi).
\]

By maximizing $\ell$ with respect to the four parameters, prove that
the maximum likelihood estimates of
$\phi,\mu_0,\mu_1,$ and $\Sigma$ are indeed as given in the formulas
above.

You may assume that there is at least one positive and one negative
example, so that the denominators in the definitions of $\mu_0$ and
$\mu_1$ above are non-zero.

% Solution to 1(d) will go here.

% ------------------------------------------------------------
% 1(e)
% ------------------------------------------------------------

\item \textbf{[3 points] Coding problem.}

In \texttt{src/p01e\_gda.py}, fill in the code to calculate
$\phi,\mu_0,\mu_1,$ and $\Sigma$, use these parameters to derive
$\theta$, and use the resulting GDA model to make predictions on the
validation set.

% ------------------------------------------------------------
% 1(f)
% ------------------------------------------------------------

\item \textbf{[5 points]}

For Dataset 1, create a plot of the training data with $x_1$ on the
horizontal axis, and $x_2$ on the vertical axis.

To visualize the two classes, use a different symbol for examples
$x^{(i)}$ with $y^{(i)}=0$ than for those with $y^{(i)}=1$.

On the same figure, plot the decision boundary found by logistic
regression in part (b).

Make an identical plot with the decision boundary found by GDA in
part (e).

% ------------------------------------------------------------
% 1(g)
% ------------------------------------------------------------

\item \textbf{[5 points]}

Repeat the steps in part (f) for Dataset 2.

On which dataset does GDA seem to perform worse than logistic
regression?

Why might this be the case?

% ------------------------------------------------------------
% 1(h)
% ------------------------------------------------------------

\item \textbf{[3 extra credit points]}

For the dataset where GDA performed worse in parts (f) and (g), can
you find a transformation of the $x^{(i)}$'s such that GDA performs
significantly better?

What is this transformation?

\end{enumerate}

% ============================================================
% Problem 2
% ============================================================


\newpage 

\item \textbf{[30 points] Incomplete, Positive-Only Labels}

In this problem we will consider training binary classifiers in
situations where we do not have full access to the labels.

In particular, we consider a scenario, which is not too infrequent in
real life, where we have labels only for a subset of the positive
examples. All the negative examples and the rest of the positive
examples are unlabeled.

That is, we assume a dataset
\[
\left\{
(x^{(i)},t^{(i)},y^{(i)})
\right\}_{i=1}^{m},
\]
where
\[
t^{(i)}\in\{0,1\}
\]
is the ``true'' label, and where
\[
y^{(i)}
=
\begin{cases}
1, & x^{(i)}\text{ is labeled},\\
0, & \text{otherwise}.
\end{cases}
\]

All labeled examples are positive, which is to say
\[
p(t^{(i)}=1\mid y^{(i)}=1)=1,
\]
but unlabeled examples may be positive or negative.

Our goal in the problem is to construct a binary classifier $h$ of
the true label $t$, with only access to the partial labels $y$.

In other words, we want to construct $h$ such that
\[
h(x^{(i)})
\approx
p(t^{(i)}=1\mid x^{(i)})
\]
as closely as possible, using only $x$ and $y$.

\textit{Real world example:}
Suppose we maintain a database of proteins which are involved in
transmitting signals across membranes. Every example added to the
database is involved in a signaling process, but there are many
proteins involved in cross-membrane signaling which are missing from
the database.

It would be useful to train a classifier to identify proteins that
should be added to the database.

In our notation, each example $x^{(i)}$ corresponds to a protein,
$y^{(i)}=1$ if the protein is in the database and $0$ otherwise, and
$t^{(i)}=1$ if the protein is involved in a cross-membrane signaling
process and thus should be added to the database, and $0$ otherwise.

\begin{enumerate}[label=(\alph*)]

% ------------------------------------------------------------
% 2(a)
% ------------------------------------------------------------

\item \textbf{[5 points]}

Suppose that each $y^{(i)}$ and $x^{(i)}$ are conditionally
independent given $t^{(i)}$:
\[
p(y^{(i)}=1\mid t^{(i)}=1,x^{(i)})
=
p(y^{(i)}=1\mid t^{(i)}=1).
\]

Note this is equivalent to saying that labeled examples were selected
uniformly at random from the set of positive examples.

Prove that the probability of an example being labeled differs by a
constant factor from the probability of an example being positive.

That is, show that
\[
p(t^{(i)}=1\mid x^{(i)})
=
\frac{
p(y^{(i)}=1\mid x^{(i)})
}{\alpha}
\]
for some
\[
\alpha\in\mathbb{R}.
\]

% Solution to 2(a) will go here.

% ------------------------------------------------------------
% 2(b)
% ------------------------------------------------------------

\item \textbf{[5 points]}

Suppose we want to estimate $\alpha$ using a trained classifier $h$
and a held-out validation set $V$.

Let $V_+$ be the set of labeled (and hence positive) examples in $V$,
given by
\[
V_+
=
\left\{
x^{(i)}\in V
\mid
y^{(i)}=1
\right\}.
\]

Assuming that
\[
h(x^{(i)})
\approx
p(y^{(i)}=1\mid x^{(i)})
\]
for all examples $x^{(i)}$, show that
\[
h(x^{(i)})
\approx
\alpha
\]
for all
\[
x^{(i)}\in V_+.
\]

You may assume that
\[
p(t^{(i)}=1\mid x^{(i)})
\approx1
\]
when
\[
x^{(i)}\in V_+.
\]

% Solution to 2(b) will go here.

% ------------------------------------------------------------
% 2(c)
% ------------------------------------------------------------

\item \textbf{[5 points] Coding problem.}

The following three problems will deal with a dataset which we have
provided in the following files:
\[
\texttt{data/ds3\_\{train,valid,test\}.csv}.
\]

Each file contains the following columns:
$x_1,x_2,y,$ and $t$.
As in Problem 1, there is one example per row.

First we will consider the ideal case, where we have access to the
true $t$-labels for training.

In \texttt{src/p02cde\_posonly.py}, write a logistic regression
classifier that uses $x_1$ and $x_2$ as input features, and train it
using the $t$-labels.

You can ignore the $y$-labels for this part.

Output the trained model's predictions on the test set to the file
specified in the code.

% ------------------------------------------------------------
% 2(d)
% ------------------------------------------------------------

\item \textbf{[5 points] Coding problem.}

We now consider the case where the $t$-labels are unavailable, so you
only have access to the $y$-labels at training time.

Add to your code in \texttt{p02cde\_posonly.py} to re-train the
classifier (still using $x_1$ and $x_2$ as input features), but using
the $y$-labels only.

% ------------------------------------------------------------
% 2(e)
% ------------------------------------------------------------

\item \textbf{[10 points] Coding problem.}

Using the validation set, estimate the constant $\alpha$ by averaging
your classifier's predictions over all labeled examples in the
validation set:
\[
\alpha
\approx
\frac{1}{|V_+|}
\sum_{x^{(i)}\in V_+}
h(x^{(i)}).
\]

Add code in \texttt{src/p02cde\_posonly.py} to rescale your
classifier's predictions from part (d) using the estimated value for
$\alpha$.

Finally, using a threshold of
\[
p(t^{(i)}=1\mid x^{(i)})=0.5,
\]
make three separate plots with the decision boundaries from parts
(c)--(e) plotted on top of the test set.

Plot $x_1$ on the horizontal axis and $x_2$ on the vertical axis, and
use two different symbols for the positive ($t^{(i)}=1$) and negative
($t^{(i)}=0$) examples.

In each plot, indicate the separating hyperplane with a red line.

\end{enumerate}

\medskip

\textit{Remark:}
We saw that the true probability $p(t\mid x)$ was only a constant
factor away from $p(y\mid x)$.

This means, if our task is to only rank examples (i.e.\ sort them) in
a particular order, then in fact we do not even need to estimate
$\alpha$.

The rank based on $p(y\mid x)$ will agree with the rank based on
$p(t\mid x)$.

% ============================================================
% Problem 3
% ============================================================


\newpage


\item \textbf{[25 points] Poisson Regression}

\begin{enumerate}[label=(\alph*)]

% ------------------------------------------------------------
% 3(a)
% ------------------------------------------------------------

\item \textbf{[5 points]}

Consider the Poisson distribution parameterized by $\lambda$:
\[
p(y;\lambda)
=
\frac{e^{-\lambda}\lambda^y}{y!}.
\]

Show that the Poisson distribution is in the exponential family, and
clearly state the values for
\[
b(y),\qquad
\eta,\qquad
T(y),\qquad
a(\eta).
\]

% Solution to 3(a) will go here.

% ------------------------------------------------------------
% 3(b)
% ------------------------------------------------------------

\item \textbf{[3 points]}

Consider performing regression using a GLM model with a Poisson
response variable.

What is the canonical response function for the family?

You may use the fact that a Poisson random variable with parameter
$\lambda$ has mean $\lambda$.

% Solution to 3(b) will go here.

% ------------------------------------------------------------
% 3(c)
% ------------------------------------------------------------

\item \textbf{[7 points]}

For a training set
\[
\left\{
(x^{(i)},y^{(i)});
\ i=1,\ldots,m
\right\},
\]
let the log-likelihood of an example be
\[
\log p(y^{(i)}\mid x^{(i)};\theta).
\]

By taking the derivative of the log-likelihood with respect to
$\theta_j$, derive the stochastic gradient ascent update rule for
learning using a GLM model with Poisson responses $y$ and the canonical
response function.

% Solution to 3(c) will go here.

% ------------------------------------------------------------
% 3(d)
% ------------------------------------------------------------

\item \textbf{[7 points] Coding problem.}

Consider a website that wants to predict its daily traffic.

The website owners have collected a dataset of past traffic to their
website, along with some features which they think are useful in
predicting the number of visitors per day.

The dataset is split into train/valid/test sets and follows the same
format as Datasets 1--3:
\[
\texttt{data/ds4\_\{train,valid\}.csv}.
\]

We will apply Poisson regression to model the number of visitors per
day.

Note that applying Poisson regression in particular assumes that the
data follows a Poisson distribution whose natural parameter is a
linear combination of the input features:
\[
\eta=\theta^Tx.
\]

In \texttt{src/p03d\_poisson.py}, implement Poisson regression for
this dataset and use gradient ascent to maximize the log-likelihood
of $\theta$.

\end{enumerate}

% ============================================================
% Problem 4
% ============================================================

\newpage

\item \textbf{[15 points] Convexity of Generalized Linear Models}

In this question we will explore and show some nice properties of
Generalized Linear Models, specifically those related to their use of
Exponential Family distributions to model the output.

Most commonly, GLMs are trained by using the negative log-likelihood
(NLL) as the loss function. This is mathematically equivalent to
Maximum Likelihood Estimation; maximizing the log-likelihood is
equivalent to minimizing the negative log-likelihood.

In this problem, our goal is to show that the NLL loss of a GLM is a
convex function with respect to the model parameters.

As a reminder, this is convenient because a convex function is one
for which any local minimum is also a global minimum.

To recap, an exponential family distribution is one whose probability
density can be represented as
\[
p(y;\eta)
=
b(y)
\exp\left(
\eta^TT(y)-a(\eta)
\right),
\]
where $\eta$ is the natural parameter of the distribution.

Moreover, in a Generalized Linear Model, $\eta$ is modeled as
\[
\theta^Tx,
\]
where
\[
x\in\mathbb{R}^n
\]
are the input features of the example, and
\[
\theta\in\mathbb{R}^n
\]
are learnable parameters.

In order to show that the NLL loss is convex for GLMs, we break down
the process into sub-parts, and approach them one at a time.

Our approach is to show that the second derivative (i.e.\ Hessian) of
the loss with respect to the model parameters is Positive
Semi-Definite (PSD) at all values of the model parameters.

We will also show some nice properties of Exponential Family
distributions as intermediate steps.

For the sake of convenience we restrict ourselves to the case where
$\eta$ is a scalar.

Assume
\[
p(Y\mid X;\theta)
\sim
\operatorname{ExponentialFamily}(\eta),
\]
where
\[
\eta\in\mathbb{R}
\]
is a scalar, and
\[
T(y)=y.
\]

This makes the exponential family representation take the form
\[
p(y;\eta)
=
b(y)\exp(\eta y-a(\eta)).
\]

\begin{enumerate}[label=(\alph*)]

% ------------------------------------------------------------
% 4(a)
% ------------------------------------------------------------

\item \textbf{[5 points]}

Derive an expression for the mean of the distribution.

Show that
\[
\mathbb{E}[Y\mid X;\theta]
\]
can be represented as the gradient of the log-partition function
$a$ with respect to the natural parameter $\eta$.

\textit{Hint:}
Start with observing that
\[
\frac{\partial}{\partial\eta}
\int p(y;\eta)\,dy
=
\int
\frac{\partial}{\partial\eta}
p(y;\eta)\,dy.
\]

% Solution to 4(a) will go here.

% ------------------------------------------------------------
% 4(b)
% ------------------------------------------------------------

\item \textbf{[5 points]}

Next, derive an expression for the variance of the distribution.

In particular, show that
\[
\operatorname{Var}(Y\mid X;\theta)
\]
can be expressed as the derivative of the mean with respect to $\eta$,
i.e.\ the second derivative of the log-partition function
$a(\eta)$ with respect to the natural parameter $\eta$.

% Solution to 4(b) will go here.

% ------------------------------------------------------------
% 4(c)
% ------------------------------------------------------------

\item \textbf{[5 points]}

Finally, write out the loss function
\[
\ell(\theta),
\]
the NLL of the distribution, as a function of $\theta$.

Then calculate the Hessian of the loss with respect to $\theta$, and
show that it is always PSD.

This concludes the proof that NLL loss of GLM is convex.

\textit{Hint:}
Use the chain rule of calculus along with the results of the previous
parts to simplify your derivations.

% Solution to 4(c) will go here.

\end{enumerate}

\medskip

\textit{Remark:}
The main takeaways from this problem are:
\begin{itemize}[leftmargin=1.5em]
    \item Any GLM model is convex in its model parameters.
    \item The exponential family of probability distributions are
    mathematically nice. Whereas calculating mean and variance of
    distributions in general involves integrals, we can calculate them
    using derivatives for exponential family distributions.
\end{itemize}

% ============================================================
% Problem 5
% ============================================================

\newpage

\item \textbf{[25 points] Locally weighted linear regression}

\begin{enumerate}[label=(\alph*)]

% ------------------------------------------------------------
% 5(a)
% ------------------------------------------------------------

\item \textbf{[10 points]}

Consider a linear regression problem in which we want to ``weight''
different training examples differently.

Specifically, suppose we want to minimize
\[
J(\theta)
=
\frac12
\sum_{i=1}^{m}
w^{(i)}
\left(
\theta^Tx^{(i)}-y^{(i)}
\right)^2.
\]

In class, we worked out what happens for the case where all the
weights (the $w^{(i)}$'s) are the same.

In this problem, we will generalize some of those ideas to the
weighted setting.

\begin{enumerate}[label=\roman*.]

% ------------------------------------------------------------
% 5(a)(i)
% ------------------------------------------------------------

\item \textbf{[2 points]}

Show that $J(\theta)$ can also be written
\[
J(\theta)
=
(X\theta-y)^T
W
(X\theta-y)
\]
for an appropriate matrix $W$, and where $X$ and $y$ are as defined
in class.

Clearly specify the value of each element of the matrix $W$.

% Solution to 5(a)(i) will go here.

% ------------------------------------------------------------
% 5(a)(ii)
% ------------------------------------------------------------

\item \textbf{[4 points]}

If all the $w^{(i)}$'s equal $1$, then we saw in class that the normal
equation is
\[
X^TX\theta=X^Ty,
\]
and that the value of $\theta$ that minimizes $J(\theta)$ is given by
\[
(X^TX)^{-1}X^Ty.
\]

By finding the derivative
\[
\nabla_\theta J(\theta)
\]
and setting that to zero, generalize the normal equation to this
weighted setting, and give the new value of $\theta$ that minimizes
$J(\theta)$ in closed form as a function of $X,W,$ and $y$.

% Solution to 5(a)(ii) will go here.

% ------------------------------------------------------------
% 5(a)(iii)
% ------------------------------------------------------------

\item \textbf{[4 points]}

Suppose we have a dataset
\[
\left\{
(x^{(i)},y^{(i)});
\ i=1,\ldots,m
\right\}
\]
of $m$ independent examples, but we model the $y^{(i)}$'s as drawn
from conditional distributions with different levels of variance
$(\sigma^{(i)})^2$.

Specifically, assume the model
\[
p(y^{(i)}\mid x^{(i)};\theta)
=
\frac{1}{\sqrt{2\pi}\sigma^{(i)}}
\exp\left(
-
\frac{
\left(
y^{(i)}-\theta^Tx^{(i)}
\right)^2
}{
2(\sigma^{(i)})^2
}
\right).
\]

That is, each $y^{(i)}$ is drawn from a Gaussian distribution with mean
\[
\theta^Tx^{(i)}
\]
and variance
\[
(\sigma^{(i)})^2,
\]
where the $\sigma^{(i)}$'s are fixed, known constants.

Show that finding the maximum likelihood estimate of $\theta$ reduces
to solving a weighted linear regression problem.

State clearly what the $w^{(i)}$'s are in terms of the
$\sigma^{(i)}$'s.

% Solution to 5(a)(iii) will go here.

\end{enumerate}

% ------------------------------------------------------------
% 5(b)
% ------------------------------------------------------------

\item \textbf{[10 points] Coding problem.}

We will now consider the following dataset, where $x^{(i)}$ is
one-dimensional:
\[
\texttt{data/ds5\_\{train,valid,test\}.csv}.
\]

In \texttt{src/p05b\_lwr.py}, implement locally weighted linear
regression using the normal equations you derived in part (a), and
using
\[
w^{(i)}
=
\exp\left(
-
\frac{
\|x^{(i)}-x\|_2^2
}{
2\tau^2
}
\right).
\]

Train your model on the train split using
\[
\tau=0.5,
\]
then run your model on the valid split and report the mean squared
error (MSE).

Finally, plot your model's predictions on the validation set. Plot the
training set with blue \texttt{x} markers and the validation set with
red \texttt{o} markers.

Does the model seem to be under- or overfitting?

% ------------------------------------------------------------
% 5(c)
% ------------------------------------------------------------

\item \textbf{[5 points] Coding problem.}

We will now tune the hyperparameter $\tau$.

In \texttt{src/p05c\_tau.py}, find the MSE value of your model on the
validation set for each of the values of $\tau$ specified in the code.

For each $\tau$, plot your model's predictions on the validation set
in the format described in part (b).

Report the value of $\tau$ which achieves the lowest MSE on the valid
split, and finally report the MSE on the test split using this
$\tau$-value.

\end{enumerate}

\end{enumerate}

\end{document}